\documentclass[11pt]{scrartcl}

\usepackage[stixtwo]{fontsetup}
\usepackage{amsmath}

\begin{document}
\begin{itemize}
	\item A számláló és a nevező legmagasabb fokú tagjával
	      ($n^{11/2}$, illetve $n^{19/3}$) egyszerűsítve:
	      \[
		      c_n
		      = \frac{n^{11/2}}{n^{19/3}} \cdot
		      \frac{1 + \frac{3}{n^{5/2}} + \frac{100}{n^{11/2}}}{3 + \frac{100}{n^{13/3}}}
		      = \frac{1}{n^{5/6}} \cdot
		      \frac{1 + \frac{3}{n^{5/2}} + \frac{100}{n^{11/2}}}{3 + \frac{100}{n^{13/3}}}
		      \to 0 \cdot \frac{1}{3} = 0
	      \]

	\item Minden $n \ge 1$ esetén
	      \[
		      \frac{n^{11/2}}{103 n^{13/3}}
		      \le \frac{n^5 \cdot \sqrt{n} + 3n^3 + 100}{3n^4 \cdot \sqrt[3]{n} + 100n^2}
		      \le \frac{104 n^{11/2}}{3 n^{13/3}},
	      \]
	      így
	      \[
		      1
		      \leftarrow \frac{\left(\sqrt[n]{n}\right)^{7/6}}{\sqrt[n]{103}}
		      \le d_n
		      \le \sqrt[n]{\frac{104}{3}} \left(\sqrt[n]{n}\right)^{7/6}
		      \to 1,
	      \]
	      tehát a rendőrelv miatt $d_n \to 1$.
\end{itemize}
\end{document}
